% Generated draft: requires mathematical and source review.
\documentclass[11pt]{article}
\usepackage[margin=25mm]{geometry}
\usepackage{fontspec}
\setmainfont{Latin Modern Roman}
\usepackage{amsmath,amssymb,mathtools,unicode-math}
\setmathfont{Latin Modern Math}
\usepackage{enumitem,needspace,xurl,hyperref}
\setlength{\parindent}{0pt}
\setlength{\parskip}{5pt}
\setlength{\emergencystretch}{3em}
\newcommand{\N}{\mathbb{N}}
\newcommand{\Z}{\mathbb{Z}}
\newcommand{\Q}{\mathbb{Q}}
\newcommand{\R}{\mathbb{R}}
\newcommand{\C}{\mathbb{C}}
\newcommand{\F}{\mathbb{F}}
\newcommand{\A}{\mathbb{A}}
\newcommand{\PP}{\mathbb{P}}
\newcommand{\E}{\mathbb{E}}
\newcommand{\eps}{\varepsilon}
\newcommand{\dd}{\,\mathrm{d}}
\begin{document}
\section*{UnsolvedMath ID: 701313}
\textbf{Source problem code:} catalogue entry
\subsection{Definitions and mathematical statement}
Let \(G\) be a group, with identity \(1\). Use the commutator convention \([x,y]=x^{-1}y^{-1}xy\). An element \(a\in G\) is a \emph{left \(3\)-Engel element} if
\[
[[[x,a],a],a]=1\qquad\text{for every }x\in G.
\]
A group is \emph{locally nilpotent} if each of its finitely generated subgroups is nilpotent. Here a group \(H\) is nilpotent if its lower central series, defined by \(\gamma_1(H)=H\) and \(\gamma_{n+1}(H)=[\gamma_n(H),H]\), reaches the trivial subgroup after finitely many steps; \([U,V]\) denotes the subgroup generated by all commutators \([u,v]\) with \(u\in U\) and \(v\in V\). The \emph{Hirsch--Plotkin radical} \(\operatorname{HP}(G)\) is the largest normal locally nilpotent subgroup of \(G\).

\textbf{Problem (A. Abdollahi, 2010).} Does there exist a group \(G\) and an element \(a\in G\) such that
\[
\bigl(\forall x\in G,\; [[[x,a],a],a]=1\bigr)
\quad\text{and}\quad
a\notin\operatorname{HP}(G)?
\]
No finiteness assumption on \(G\) is imposed in the supplied statement.
\subsection{Short English statement}
Does some group contain a left \(3\)-Engel element outside its Hirsch--Plotkin radical, that is, outside its largest normal locally nilpotent subgroup?
\subsection{Sources}
\begin{itemize}
\item Local UnsolvedMath registry: permanent ID \(701313\), source problem code \texttt{catalogue entry}, title ``Kourovka Notebook Problem 17.11''; supplied upstream revision \texttt{d371ee917ac9647fbb06cd6fc529deec6f3dda82}.
\item The supplied background attributes Problem 17.11 to A. Abdollahi, proposed in 2010, in \emph{The Kourovka Notebook (Unsolved Problems in Group Theory)}, 17th issue (2010), and cites the 21st issue edition (2026), \texttt{arXiv:1401.0300v48}. The supplied collection description names E. I. Khukhro and V. D. Mazurov as editors.
\item The registry labels the problem \texttt{open}. This is a supplied source classification, not an independently verified assessment of the literature. The cited editions and arXiv version have not been independently checked for this draft.
\end{itemize}
\end{document}
